\section{Related work}\label{sec:related-work}

The closest comparison is with \citet[][Theorems 1--2 and Remarks 2 and 11]{KennedyBalakrishnanRobinsWasserman2024}, who characterize pointwise conditional average treatment effect estimation under separate H\"older restrictions on the propensity score, control regression, and treatment contrast. For their formula, dimension one, contrast smoothness one, and average nuisance smoothness one tenth give the rate \(n^{-1/4}\); one tenth is below the corresponding one-sixth threshold. The present choice therefore isolates a simple rough-nuisance regime in which higher-order bias control is relevant. Their lower bound applies to a broader experiment allowing real outcomes; its construction uses a binary-outcome submodel and a covariate density bounded above that can vanish locally. Their higher-order local-polynomial estimator attains the stated rate under additional approximation, eigenvalue, boundedness, and covariate-distribution estimation conditions. Remark 11 expresses the last requirement through a third-order product involving design estimation and nuisance errors. The present experiment in \cref{obj:def:complete-model} combines binary outcomes, a scalar covariate, a density bounded above and away from zero, fixed overlap, one-tenth H\"older nuisance functions, and a Lipschitz treatment contrast. Within this class, \cref{obj:thm:matched-risk-frontier} gives matching orders for private pointwise expected absolute error, with additional scales \((n^2\epsilon)^{-1/7}\) and \((n\epsilon)^{-1/2}\). The comparison therefore concerns both the statistical experiment and the attainment argument: here, occupancy-weighted within-cell pairs support private attainment under the stated density bounds.
% Source: [Minimax rates for heterogeneous causal effect estimation](https://arxiv.org/pdf/2203.00837).

This construction connects to the literature on private U-statistics. \citet[][Theorem 2]{ChaudhuriLohPandeySarkar2024} give a pure differentially private estimator for the expectation of a bounded symmetric statistic, with a probability bound involving its variance, concentration of local H\'ajek projections, and privacy terms. The release in \cref{obj:def:private-ratio} combines two count-weighted pair statistics within a public partition, adds Laplace noise jointly, and forms a clipped ratio with a public denominator floor. Its analysis addresses the random number of observations contributing pairs and the normalization needed for a local treatment contrast. The shared occupancy weights in the population numerator and denominator are central to \cref{obj:lem:occupancy-cancellation}, while \cref{obj:lem:all-count-stability} controls replacement sensitivity across cell counts, including transitions involving cells with fewer than two observations. Together, these features support the uniform guarantee in \cref{obj:thm:uniform-private-upper}.
% Source: [On differentially private U statistics](https://proceedings.neurips.cc/paper_files/paper/2024/file/290848c892a85a6936b473d4daa6d0ad-Paper-Conference.pdf).

The lower-bound analysis draws on coupling-based private testing. \citet[][Section 2]{AcharyaSunZhang2021} relate private distinguishability to the expected Hamming distance between datasets under a coupling, including comparisons between mixtures of product laws. The present analysis specializes this approach to causal alternatives satisfying the binary experiment and its smoothness restrictions. The common-mass construction in \cref{obj:def:common-mass-coupling} operates on conditional components of the shared-sign experiment, and \cref{obj:lem:measurable-component-contraction,obj:lem:private-kernel-comparison} connect these comparisons to private output laws. The separated testing priors in \cref{obj:lem:frontier-testing-priors} then supply the experiments used for the estimation and interval lower bounds. The contribution lies in constructing and controlling these comparisons within the specified causal class.
% Source: [Differentially private Assouad, Fano, and Le Cam](https://proceedings.mlr.press/v132/acharya21a/acharya21a.pdf).

Private heterogeneous-effect learners provide a complementary perspective. \citet{NiuEtAl2022} develop a general procedure for conditional average treatment effect learning under approximate differential privacy, using sample splitting and private base learners. Their implementations include single-stage and multistage estimators, and their empirical performance criterion is test-set mean squared error, with a bias--variance decomposition. \citet[][Theorem 1]{SchroderMelnychukFeuerriegel2025} instead establish approximate differential privacy for a prespecified finite vector of conditional treatment effect queries through Gaussian output perturbation calibrated to the second-stage learner's sensitivity. Their accompanying performance analysis concerns Neyman orthogonality and quasi-oracle estimation under nuisance estimation conditions \citep[][Theorem 2]{SchroderMelnychukFeuerriegel2025}. The present estimand is the treatment contrast at a fixed covariate location, privacy is pure differential privacy under replacement adjacency as defined in \cref{obj:def:private-kernels}, and performance is measured by maximal expected absolute error and honest maximal expected interval length in \cref{obj:def:risk-criterion,obj:def:interval-criterion}. These choices make the comparisons specific to the estimand, privacy requirement, and decision criterion.
% Sources: [Differentially private estimation of heterogeneous causal effects](https://proceedings.mlr.press/v177/niu22a/niu22a.pdf), [Differentially private learners for heterogeneous treatment effects](https://arxiv.org/html/2503.03486v1).

The methodological background includes residualization in semiparametric regression \citep{Robinson1988}, orthogonal scores and cross-fitting \citep{Chernozhukov2018}, and two-stage conditional treatment effect estimation \citep{Nie2021,Kennedy2023}. Higher-order influence functions supply a general framework for controlling nonlinear nuisance bias \citep{Robins2008}. Causal forests and generalized random forests provide adaptive local estimation and asymptotic inference under their respective regularity conditions \citep{WagerAthey2018,Athey2019}. These approaches explain the importance of separating contrast regularity from nuisance complexity; the present pair construction gives a concrete private implementation for the rough-nuisance experiment in \cref{obj:def:complete-model}.
% Sources: [Root-N-consistent semiparametric regression](https://doi.org/10.2307/1912705), [Double/debiased machine learning for treatment and structural parameters](https://doi.org/10.1111/ectj.12097), [Quasi-oracle estimation of heterogeneous treatment effects](https://arxiv.org/abs/1712.04912), [Towards optimal doubly robust estimation of heterogeneous causal effects](https://arxiv.org/abs/2004.14497), [Higher order influence functions and minimax estimation of nonlinear functionals](https://arxiv.org/abs/0805.3040), [Estimation and inference of heterogeneous treatment effects using random forests](https://arxiv.org/abs/1510.04342), [Generalized random forests](https://imstat.org/publications/aos/aos_47_2/aos_47_2.pdf).

Smooth-nuisance kernel results form a separate model comparison. \citet[][Sections 2, 5, and 6]{Kim2026} study conditional treatment effects under reproducing kernel Hilbert space restrictions on response functions, with additional spectral, source, or dimension structure for the contrast. Their kernel ridge regression analysis relates estimation rates to contrast complexity and overlap within those models. Here, regularity is specified directly through the H\"older and Lipschitz conditions in \cref{obj:def:complete-model}, and the private risk comparison is established for that class. The distinction between these function spaces and performance criteria governs the interpretation of their respective optimality results.
% Source: [Optimal and structure-adaptive CATE estimation with kernel ridge regression](https://arxiv.org/html/2602.18958v1).

Finally, the interval criterion connects private causal inference to honest nonparametric inference. \citet{ArmstrongKolesar2018} derive finite-sample optimal intervals and efficiency bounds for linear regression functionals under Gaussian errors with known variance, while \citet{ArmstrongKolesar2020} develop intervals whose critical values account explicitly for worst-case bias. The present construction likewise incorporates a public error envelope, and \cref{obj:thm:sharp-interval-frontier} establishes matching orders, up to constant factors, for maximal expected length subject to uniform finite-sample \(90\%\) coverage over the causal class and pure replacement privacy. A complementary private causal inference result is provided by \citet[][Sections 4.2--4.3]{SchroderHartensteinFeuerriegel2025}, who construct approximate differentially private intervals for the average treatment effect using asymptotic normality and accounting for uncertainty from estimation and privatization. The present expected-length result places pointwise inference for a Lipschitz contrast with rough nuisance functions within a uniform finite-sample decision framework.
% Sources: [Optimal inference in a class of regression models](https://www.princeton.edu/~mkolesar/papers/optimal.pdf), [Simple and honest confidence intervals in nonparametric regression](https://onlinelibrary.wiley.com/doi/10.3982/QE1199), [PrivATE: Differentially private confidence intervals for average treatment effects](https://arxiv.org/html/2505.21641v1).
